% TOP_PROBLEM: {"problemId": "problem.mumford-shah-conjecture-on-the-jump-set-of-planar-minimizers", "canonicalTitle": "Planar homogeneous Mumford–Shah local-structure conjecture", "releaseRank": 279, "primaryDomain": "analysis_pde", "primaryDomainLabel": "Analysis & PDE", "displayStatus": "open_with_solved_subcases", "releaseStatus": "open_with_solved_subcases"}
% !TeX program = lualatex
% Definition and sources only; this file is not an LLM solution attempt.
\documentclass[11pt]{article}
\usepackage[margin=25mm]{geometry}
\usepackage{fontspec}
\setmainfont{Latin Modern Roman}
\usepackage{amsmath,amssymb,mathtools}
\usepackage{unicode-math}
\setmathfont{Latin Modern Math}
\usepackage{xurl,hyperref}
\hypersetup{colorlinks=true,urlcolor=blue}
\setcounter{secnumdepth}{0}
\setlength{\parindent}{0pt}
\setlength{\parskip}{5pt}
\setlength{\emergencystretch}{3em}
\begin{document}
\section*{\texorpdfstring{Planar homogeneous Mumford–Shah local-structure conjecture}{Planar homogeneous Mumford–Shah local-structure conjecture}}\label{279}
\subsection{Definitions and mathematical statement}
For $u\in W^{1,2}_{\mathrm{loc}}(\Omega\setminus K)$ and relatively closed countably rectifiable $K\subset\Omega\subset{\mathbb{R}}^2$, the homogeneous Mumford--Shah energy is
$E_0(u,K;U)=\int_{U\setminus K}|\nabla u|^2+\mathcal H^1(K\cap U)$.
Use reducedness and absolute compactly supported local minimality exactly as stated below, allowing all admissible changes of the crack set, without a separation constraint. The conjecture says that near each $x\in K$, a $C^1$ diffeomorphism fixing $x$ with derivative $I$ straightens $K$ to a diameter, a single radius ending at $x$, or three radii meeting at $120$ degrees. Thus only a regular crack, crack tip, or triple junction is allowed. The source's hypotheses on local finite energy and negligible removable crack pieces are part of the statement; this is not a classification of arbitrary stationary pairs.
\par\smallskip\textit{Formulation and concept references:} \href{https://arxiv.org/pdf/2308.14660v2}{[S1]}.

\subsection{Short English statement}
Must every reduced planar Mumford–Shah local minimizer have only smooth crack arcs, crack tips, and three-way junctions with 120-degree angles?
\subsection{Sources}
\begin{itemize}
\item[S1]The regularity theory for the Mumford-Shah functional on the plane; arXiv2308.14660v2 January5,2025.
\url{https://arxiv.org/pdf/2308.14660v2}.
\textit{Formulation source.}
\end{itemize}
\end{document}
